Chemistry-conditioned generative models are increasingly used to simulate how a named intervention would change physiological features. Signed concordance, continuous effect-vector fidelity, magnitude transfer, and residual encoding must be scored separately; conflating them invites failure modes in which one metric succeeds while another fails. Conditional variational autoencoders (CVAEs) make that risk concrete, because reconstruction, perturbation fidelity, and downstream classification need not move together \citep{Kingma2014VAE,Higgins2017betaVAE,Rampasek2019DrVAE,Lotfollahi2019scGen}. A formal evaluation design that measures those quantities on the same EEG-drug generator remains scarce.

Closest methods address parts of the problem: standard and \(\beta\)-VAEs \citep{Kingma2014VAE,Higgins2017betaVAE}, molecular or single-cell perturbation models \citep{Rampasek2019DrVAE,Lotfollahi2019scGen}, and diagnosis-oriented EEG twins such as DADD \citep{Amato2025DADD,Amato2026DADDnpj}. Alzheimer's disease (AD) drug response is a natural stress test given heterogeneous AChEI benefit and established EEG markers of cholinergic change \citep{Lu2020DonepezilGenetics,Pozzi2022AChEIPredictors,Jeong2004EEGreview,Cassani2018EEGADReview,Babiloni2006DonepezilEEG}. Those facts motivate constrained in silico simulation of Donepezil and Memantine. They do not license treating every latent as a clinical biomarker or every signature match as post-dose validation.

This paper contributes \textbf{SEME} (Signed concordance, Effect-vector fidelity, Magnitude transfer, Encoding retention): a formal evaluation protocol for literature-constrained, chemistry-conditioned EEG generators. The worked example is a pharmacodynamically penalized CVAE with 2185-D resting EEG features and frozen ChemBERTa embeddings for Donepezil and Memantine \citep{Chithrananda2020ChemBERTa,Winblad2007Memantine}. Under SEME we score signed concordance and continuous effect-vector correlation on Temple University Hospital (TUH) EEG \citep{Obeid2016TUH}, an Open Science Framework (OSF) AD/healthy control set, and P-ADIC AD/healthy control recordings \citep{Shor2021PADIC}; magnitude diagnosis across an escalating sequence that adds CAUEEG Normal/MCI/Dementia strata \citep{Kim2023CAUEEG}; and CAUEEG encoding probes against spectral and raw-feature references. Training features use OpenNeuro ds004504 \citep{Miltiadous2023OpenNeuro}. Domain shift is part of the test \citep{Gulrajani2021DomainShift}.

The paper does not claim clinical validation, pharmacological efficacy, post-dose digital-twin validation, diagnostic superiority, ChemBERTa chemical generalization, or constraint-caused diagnostic compression. The contributions are:

\begin{itemize}
  \item SEME as a decoupled protocol with matched unconstrained, prior-only (circularity budget), linear imitation, five-fold, and multi-seed controls as core method.
  \item Constrained twins exceed unconstrained twins on the locked sign endpoint (10/10 versus 3/10--4/10 on TUH/OSF/P-ADIC), while continuous \(r\) at locked fold-0 seed-42 (\(0.908/0.851/0.918\)) is fold- and seed-sensitive (five-fold means \(0.385/0.371/0.372\); non-42 seeds typically 0.08 to 0.43) and prior-only already reaches 10/10 signs at mean \(r = 0.636\).
  \item Magnitude does not diagnose reliably across four escalating external tests, and an encoding bake-off retains only partial label information in \(\mu_{\mathrm{base}}\) (latent ROC-AUC 0.579; best head OOF 0.593) below theta/alpha 0.675 and a 2185-D logistic feature-space reference of 0.699.
\end{itemize}

Section 2 defines the case-study twin and SEME endpoints. Section 3 applies the protocol. Section 4 interprets incremental value and dissociation. Section 5 concludes.
